\documentclass[10pt]{article}
\usepackage{foundation-backbone}

\begin{document}
\foundationtitle{Foundation as an Executable Application Backbone}{A five-page technical guide for senior developers across PostgreSQL, JavaScript, mobile, edge and native ecosystems}{Audience: senior developers and architects who may be new to Clojure, Lisp and Foundation}

\section{1. Architectural model}

Foundation Base is not primarily a web framework. It is a language and runtime toolkit for describing behaviour, transforming source forms, emitting target-language code, executing that code in target runtimes and testing agreement across those stages.

For application work, three areas form the backbone:

\begin{itemize}
  \item \product{postgres.core}: PostgreSQL-oriented definitions and database-side RPC operations;
  \item \product{xt.substrate}: nodes, spaces, request transport, eventing and runtime context;
  \item \product{xt.db.node}: models, named views, local stores, invalidation and remote refresh.
\end{itemize}

For developers unfamiliar with Lisp, a form is both executable code and inspectable data. That permits Foundation to transform the same logical definition into supported target syntax.

\begin{center}
\begin{tikzpicture}[node distance=9mm and 12mm,scale=.86,transform shape]
  \node[fbwide] (form) {Source form\\data structure};
  \node[fbwide,right=of form] (prep) {Preprocessing\\and macros};
  \node[fbwide,right=of prep] (grammar) {Target grammar\\supported forms};
  \node[fbwide,below=of prep] (emit) {Emitter output};
  \node[fbwide,left=of emit] (runtime) {Target runtime};
  \node[fbwide,right=of emit] (tests) {Emitter and\\runtime tests};
  \draw[fbedge] (form) -- (prep);
  \draw[fbedge] (prep) -- (grammar);
  \draw[fbedge] (grammar) -- (emit);
  \draw[fbedge] (emit) -- (runtime);
  \draw[fbedge] (emit) -- (tests);
  \draw[fbedge] (runtime) -- (tests);
\end{tikzpicture}
\end{center}

A portable module is therefore more than a shared library. Its source form, preprocessing, grammar representation, emitted syntax and runtime result all belong to the implementation contract.

\callout{Technical proposition}{Treat the application model and its data movement as executable infrastructure. Generate thin target adapters only after RPC, refresh and projection behaviour are proven.}

\newpage
\section{2. Spaces, views and PostgreSQL}

A substrate node owns application spaces. A space represents an execution context such as \texttt{room/local} or \texttt{room/server}. \product{xt.db.node} installs models into spaces. A model groups named views; a view holds query input, status, value, dependencies and a query key.

The repository walkthroughs exercise an explicit local SQLite and remote PostgreSQL pipeline:

\begin{center}
\begin{tikzpicture}[node distance=9mm and 14mm,scale=.83,transform shape]
  \node[fbwide] (ui) {UI adapter};
  \node[fbwide,right=of ui] (local) {local live view};
  \node[fbstorage,below=of local] (sqlite) {SQLite};
  \node[fbwide,right=of local] (transport) {substrate transport};
  \node[fbwide,right=of transport] (server) {server model / RPC};
  \node[fbstorage,below=of server] (pg) {Postgres};
  \draw[fbedge] (ui) -- node[above,scale=.75]{refresh/read} (local);
  \draw[fbedge] (local) -- (transport);
  \draw[fbedge] (transport) -- (server);
  \draw[fbedge] (server) -- (pg);
  \draw[fbedge] (local) -- (sqlite);
  \draw[fbdashed,bend right=27] (server.south) to node[below,scale=.75]{project result} (local.south);
\end{tikzpicture}
\end{center}

The application-facing model is defined before storage. A view may begin in an in-memory cache, move to SQLite or delegate to a remote Postgres-backed space without changing the conceptual screen contract.

\subsection{Role of \product{postgres.core}}

\product{postgres.core} is part of the target-language system rather than an ORM over SQL. PostgreSQL functions and supporting definitions can be represented in Foundation forms, emitted as SQL/PLpgSQL and executed against PostgreSQL.

\begin{center}
\begin{tikzpicture}[node distance=7mm,scale=.86,transform shape]
  \node[fbwide] (src) {Foundation Postgres form};
  \node[fbwide,below=of src] (sql) {Emitted SQL / PLpgSQL};
  \node[fbwide,below=of sql] (db) {Database execution};
  \node[fbwide,below=of db] (rpc) {RPC invocation};
  \node[fbwide,below=of rpc] (projection) {Local view projection};
  \draw[fbedge] (src) -- (sql);
  \draw[fbedge] (sql) -- (db);
  \draw[fbedge] (db) -- (rpc);
  \draw[fbedge] (rpc) -- (projection);
\end{tikzpicture}
\end{center}

Validation should separately prove emitted syntax, database execution, RPC behaviour and the resulting local projection. A passing emitter test alone does not establish runtime correctness.

\newpage
\section{3. Testing and AI-generated adapters}

The highest-value tests prove agreement across layers rather than only isolated CRUD functions.

\begin{center}
\begin{tikzpicture}
  \filldraw[fill=FoundationLight,draw=FoundationBlue] (0,0) -- (12,0) -- (10.8,1.25) -- (1.2,1.25) -- cycle;
  \node at (6,.62) {grammar, transformation and emitter facts};
  \filldraw[fill=FoundationLight,draw=FoundationTeal] (1.2,1.45) -- (10.8,1.45) -- (9.5,2.7) -- (2.5,2.7) -- cycle;
  \node at (6,2.07) {target runtime and database execution};
  \filldraw[fill=FoundationLight,draw=FoundationGold] (2.5,2.9) -- (9.5,2.9) -- (8.2,4.15) -- (3.8,4.15) -- cycle;
  \node at (6,3.52) {RPC, sync and projection flows};
  \filldraw[fill=FoundationLight,draw=FoundationRed] (3.8,4.35) -- (8.2,4.35) -- (7.1,5.6) -- (4.9,5.6) -- cycle;
  \node[align=center] at (6,4.97) {thin UI acceptance};
\end{tikzpicture}
\end{center}

A representative flow test should:

\begin{enumerate}
  \item create controlled server and local state;
  \item invoke an RPC or remote view refresh;
  \item assert the authoritative PostgreSQL result;
  \item assert local SQLite projection, status and invalidation;
  \item repeat for permission and failure paths.
\end{enumerate}

Once that backbone is executable, AI-generated UI becomes adapter generation rather than domain invention.

\begin{center}
\begin{tikzpicture}[node distance=9mm and 14mm,scale=.87,transform shape]
  \node[fbwide] (backbone) {Executable backbone};
  \node[fbwide,below=of backbone] (adapter) {Generated adapter boundary};
  \node[fbbox,below left=of adapter] (next) {Next.js};
  \node[fbbox,below=of adapter] (expo) {Expo / RN};
  \node[fbbox,below right=of adapter] (edge) {Bun / edge};
  \draw[fbedge] (backbone) -- (adapter);
  \draw[fbedge] (adapter) -- (next);
  \draw[fbedge] (adapter) -- (expo);
  \draw[fbedge] (adapter) -- (edge);
\end{tikzpicture}
\end{center}

The AI should receive view definitions, statuses, result shapes, RPC arguments, permission outcomes, fixtures, passing flow tests and target-framework constraints. The generated client renders state and invokes operations; it should not reproduce database rules.

This arrangement keeps UI acceptance tests thin. They verify interaction, accessibility and target integration while emitter, runtime, RPC and sync tests carry most correctness.

\newpage
\section{4. Ecosystem comparison and rewrite strategy}

\begin{tabularx}{\textwidth}{P{27mm}YYY}
\toprule
\textbf{Approach} & \textbf{Contract} & \textbf{Local state} & \textbf{Main trade-off} \\
\midrule
REST/OpenAPI & HTTP operations and schemas & Client responsibility & Strong stubs; weak view lifecycle model \\
GraphQL & Type graph and resolvers & Framework-specific cache & Flexible queries; distributed mutation semantics \\
gRPC/Protobuf & Binary service schema & Application responsibility & Excellent stubs; RPC-centric \\
ORM/Prisma & Database model and client & Application responsibility & Strong DB access; domain and sync remain elsewhere \\
Supabase/Firebase & Vendor schema, auth and SDK & Platform-dependent & Rapid delivery; vendor semantics dominate \\
CQRS/event sourcing & Commands, events, projections & Explicit & Strong history; high operational complexity \\
Foundation & Forms, grammars, spaces, views and RPC & Explicit local projection & Cross-stage discipline and learning curve \\
\bottomrule
\end{tabularx}

\subsection{Foundation versus a Bun/Rust rewrite}

A Bun/Rust rewrite is a runtime-first optimisation. Foundation is a semantics-first optimisation.

\begin{tabularx}{\textwidth}{P{31mm}YY}
\toprule
\textbf{Dimension} & \textbf{Foundation} & \textbf{Bun/Rust rewrite} \\
\midrule
Source of truth & Portable forms plus emitter/runtime tests & Protocol schemas plus Rust/TS implementations \\
Primary benefit & Cross-target consistency and generation & Native performance and direct runtime control \\
Local/remote sync & Modelled through spaces, views and projections & Requires a selected sync layer or custom protocol \\
Debugging & Trace source form through emitted target & Target-native debugging per service \\
Change surface & One model may regenerate several targets & Contract changes span services and SDKs \\
Best fit & Cross-cutting, sync-heavy applications & Measured CPU, memory, latency or startup constraints \\
\bottomrule
\end{tabularx}

A hybrid is often stronger than a repository-wide rewrite: retain Foundation for contract generation and projection tests; use Bun as a TypeScript-native web or edge adapter; use Rust for a narrow, profiled service boundary.

\subsection{Implementation decision rules}

\begin{itemize}
  \item Prefer Foundation when the same behaviour must remain aligned across PostgreSQL, web, mobile and edge targets.
  \item Prefer direct Bun/Rust when profiling proves that runtime characteristics materially constrain the product.
  \item Keep native services behind stable RPC boundaries with conformance tests.
  \item Verify emitted SQL separately from execution and execution separately from sync.
  \item Follow adjacent Foundation source, grammars and \product{code.test} facts before introducing abstractions.
\end{itemize}

\callout{Practical rule}{Optimise semantic duplication first. Optimise runtime hotspots second. Preserve an executable contract across both.}

\end{document}
